\title{X-LoRA: Mixture of Low-Rank Adapter Experts, a Flexible Framework for Large Language Models with Applications in Protein Mechanics and Molecular Design}

\begin{abstract}
We present a mixture of expert approach to develop fine-tuned large language models by employing a deep, layer-wise, token-level method based on low-rank adaptation (LoRA). Beginning with a collection of pre-trained LoRA adapters, our gating mechanism leverages hidden states to dynamically combine adapted layers, enabling the resulting X-LoRA model to harness diverse capabilities and generate novel deep layer-wise combinations to address various tasks. This design draws inspiration from biological principles of universality and diversity, where neural network components are reutilized in multiple hierarchical forms. Consequently, the X-LoRA framework can be seamlessly integrated with any existing large language model (LLM) without requiring modifications to the original architecture. We develop a specialized X-LoRA variant that delivers scientific functionalities including forward and inverse analysis tasks alongside enhanced reasoning abilities, concentrating on biomaterial analysis, protein mechanics, and design. The significance of this work lies in providing accessible, scalable, and adaptable models enriched with robust domain expertise and the ability to unify knowledge from different disciplines. Incorporating experts in biology, mathematics, reasoning, bio-inspired materials, mechanics, chemistry, protein biophysics, mechanics, and quantum mechanics-based molecular properties, we conduct a series of physics-centered case studies. These investigations cover knowledge recall, protein mechanics forward and inverse tasks, protein design, adversarial agentic modeling including ontological knowledge graph creation, and molecular design. The model not only predicts quantitative nanomechanical properties of proteins and quantum mechanical molecular attributes but also reasons about these results and accurately infers probable mechanisms underlying distinct molecular phenomena.
\end{abstract}

\section{Introduction}
Large language models (LLMs)~\cite{Vaswani2017AttentionNeedc,Touvron2023LlamaModels,OpenAI2023GPT-4Report,Chowdhery2022PaLM:Pathways,Jiang2023Mistral7B,Gunasekar2023TextbooksNeed} have become widely popular, including the creation of specialized models that excel in specific tasks, reasoning processes, or scientific fields~\cite{Bubeck2023SparksGPT-4,Buehler2023MechGPTModalities,Nejjar2023LLMsAnalysis,Buehler2023GenerativeDesign,Luu2023BioinspiredLLM:Materials,Luu2023GenerativeSolvents,Buehler2023MeLMProblemsb,Ge2023OpenAGI:Experts}. However, training these models can be expensive, particularly when a broad range of capabilities is required. Approaches like low-rank adapters (LoRA) \cite{hu2021lora} have been introduced as a more resource-efficient alternative, though they typically target more specialized knowledge domains. The core idea behind LoRA is to incorporate low-rank matrices that augment the original full-sized matrix, designating these low-rank matrices as the sole trainable parts of the model. Because only the adapter layers undergo training, these models generally retain the pre-trained knowledge from the base model while adapting the model to particular tasks, all while maintaining computational efficiency. Due to the efficiency of training LoRA adapters, it becomes feasible to create a large number of specialized models with this method.

Although LoRA adapters offer an effective means to develop specialized capabilities, the challenge of how to combine multiple such systems into enhanced models that provide an integrated and potentially improved range of functionalities remains unresolved. Indeed, combining LLMs has emerged as a compelling research topic~\cite{Kim2023SOLARUp-Scaling}. Additionally, experimental techniques such as spherically interpolated model parameters (SLERP) and related methods~\cite{Arcee-ai/mergekit:Models.} have yielded promising results, demonstrating the potential of merging models to unlock novel capabilities. Models combined in these manners are typically generated {\it a priori} using specific algorithms; while they have shown encouraging performance on certain benchmarks, further investigation is needed to develop principled approaches for constructing mixed models and to understand how particular features manifest.

We propose that dynamically mixing parameters may represent a more generalizable method, as it enables the mixed model to investigate novel combinations during inference.

Similar methodologies have employed mixture of experts (MoE) frameworks~\cite{Jacobs1991AdaptiveExperts,Hampshire1992TheRecognition,Jordan1994HierarchicalAlgorithm}, including the recently introduced Mixtral LLM~\cite{Jiang2024MixtralExperts}. Nonetheless, these methods often incur high computational costs and substantial resource demands even at inference time. In this work, we introduce a computationally efficient mixture-of-expert technique that utilizes a collection of distinct LoRA adapters, each of which can be trained independently with ease. Our strategy adopts a fine-grained and adaptable approach, drawing inspiration from a biological model that repurposes neural network components to build a hierarchy of interacting models, ranging from the foundational model up to agentic systems (Fig.~\ref{fig:hierarchical_concept}). To quantify its effectiveness, our study specifically applies this strategy to science-oriented LLMs, focusing in particular on protein mechanics challenges within the broader domain of bio-inspired materials analysis and design.

We present an algorithm for dynamically mixing adapters at the token-level state, with a high degree of granularity that permits adapter mixing at the scale of individual layers. This method, termed X-LoRA, provides access to a variety of capabilities crucial for scientific applications of multimodal language modeling~\cite{Hu2023DeepScience,Buehler2023AMetamaterials,Buehler2022ProteinNetworks,}. 

\subsection{Fundamental concepts of X-LoRA}

We begin with a brief overview of the principles underlying the development of LoRA adapters. The core idea behind low-rank adaptation (LoRA)~\cite{hu2021lora} posits that weight updates occupy a low “intrinsic dimension,” leveraging this by keeping the original weights $W_0 \in \mathbb{R}^{d \times k}$ fixed and restricting the updates to a low-rank decomposition

\begin{equation}
W_0 + \Delta W_0 = W_0 + BA 
\end{equation}

where $B \in \mathbb{R}^{d \times r}$ and $A \in \mathbb{R}^{r \times d}$, with the rank $r$ much smaller than ${\rm min}(d,k)$. In the training process using LoRA~\cite{hu2021lora}, the pretrained weights $W_0$ remain fixed, and only the matrices $A$ and $B$ are updated as trainable parameters. The forward propagation of a LoRA layer can be described by the equation:

\begin{equation}
    h = W_0x + \Delta W_0x = W_0x + BAx
\end{equation}

It should be noted that, practically, LoRA adapters include a fixed scalar parameter $\alpha$, which is applied to the factorized matrices, resulting in:

\begin{equation}
    h = W_0x + BAx \cdot \alpha
\end{equation}

Building upon this concept, X-LoRA introduces the idea of employing multiple LoRA adapters, each endowed with distinct capabilities, as demonstrated in prior scientific LLM research (e.g., \cite{Buehler2023GenerativeDesign}). Instead of statically mixing or merging models, we propose a dynamic gating mechanism that adjusts individual LoRA adapters with granularity at the token and layer levels, allowing mixing deeper within the model. More specifically, for a given LoRA adapter $i$, its original $\alpha_i$ parameter is scaled by a factor denoted $\lambda_i$, producing a new scaling parameter for the adapter, $\alpha^*_i$:

\begin{equation}
    \alpha^*_i = \alpha_i \cdot \lambda_i
\end{equation}

The scaling factor $\lambda_i$ is generated by an X-LoRA scaling head which leverages the model’s hidden states, constituting the trainable component in the X-LoRA architecture. Because the scaling head can utilize complex representations, it can be efficiently implemented as a straightforward feed-forward neural network comprising one or a few layers (in the X-LoRA implementation, users can specify the neural network's design details).

The scaling is calculated for each token in the input sequence, thereby offering a fine-grained control since each LoRA adapter operating at particular layers of the LLM is scaled individually.

From this point onward, we refer to this technique of modulating the components of the adapters in this manner as scaling. Additional information regarding this method is provided in the Materials and Methods section.

\subsection{Paper outline}

The structure of this paper is organized as follows. Initially, we describe the methodology and training approach, including the relevant datasets that contribute to the creation of an X-LoRA model with specialized abilities in the physical sciences, with a particular emphasis on biomaterials such as proteins. Next, we demonstrate a series of experiments where we apply the X-LoRA to a range of tasks, including question answering, conversational and agentic modeling, protein design and analysis, among others. We provide a thorough examination of the scaling behaviors and validate our method by comparing it to molecular modeling and other physical datasets and techniques.

\section{Results and discussion}

The development of the X-LoRA model involves the following stages:

\begin{enumerate}
    \item Training a foundational base LLM (completed in prior work) \item Training individual adapters separately to gain expertise in various domains (these represent distinct or overlapping skill and knowledge areas) \item Training the integrated X-LoRA model
\end{enumerate}

Our experimental process begins with the training of multiple LoRA adapters. We create nine adapters, each fine-tuned for specific expertise, based on the Zephyr-7B-$\beta$ model~\cite{HuggingFaceH4/zephyr-7b-betaFace}, which itself was developed atop the Mistral-7B model\cite{Jiang2023Mistral7B}:

\begin{enumerate}
    \item Bioinspired materials \item Chain-of-thought (CoT) and reasoning \item Chemistry \item Mathematics \item Physics \item Biology \item Mechanics and materials \item Platypus: Logic and reasoning \item Protein mechanics tasks (including generative sequence-to-property and inverse modeling capabilities)
\end{enumerate}

These specific domains of concentration and expertise were selected because our ultimate goal is to develop a STEM-oriented LLM with a focus on materials science. Nevertheless, we hypothesize that our methodology can be broadly extended to various fields, particularly to generate modalities that facilitate connections between different knowledge areas~\cite{Buehler2023MechGPTModalities,Buehler2023MeLMProblemsb}.

Once the LoRA adapters are established, we proceed to train the X-LoRA model. For this, we utilize a subset derived from the combined pool of training samples that were employed to train the individual LoRA adapters (further details are provided in the Materials and Methods section, including comprehensive descriptions of the dataset, implementation, and related aspects).

We treat the X-LoRA system as an experimental framework to evaluate its performance across several scenarios. We carry out multiple experiments, such as tackling complex question answering, employing the generative model for protein design, conducting analysis and reasoning, and developing more advanced interaction techniques through adversarial agentic modeling. Throughout these evaluations, we compare the performance of the base model against that of the X-LoRA model, and monitor how the scaling head leverages various layers from different experts during the generation or formulation of responses.

\subsection{Question answering and observed X-LoRA layer-wise scaling weights}

We begin by exploring the distinctions between the X-LoRA model and the foundational base model. We pose the following question to both models:

\begin{quote}
\small\texttt{What material behavior (stiffening vs softening) causes supersonic crack speeds in mode II cracks? Think step by step.}
\end{quote}

The base model replies:

\begin{quote}
\small\texttt{Supersonic crack speeds, also known as Rayleigh waves, occur in mode II cracks due to the stiffening material behavior of the material near the crack tip. In mode II fracture, the loading is shear-dominated and the crack propagates perpendicular to the direction of applied stress. Here's a step-by-step explanation:

As the crack tip nears the material, the stresses around it increase sharply. This leads to a zone of elevated stress concentrations known as the stress singularity.

Within this stress singularity, the material undergoes elastic deformation, resulting in enhanced stiffness, i.e., a stiffening effect. This stiffening occurs because the material adjacent to the crack tip becomes more restricted as the crack advances [...]}
\end{quote}

Based on our domain expertise~\cite{Buehler2008AtomisticFailure}, this response is not only inaccurate but also overly verbose. In contrast, the X-LoRA model provides the following reply:

\begin{quote}
\small\texttt{The material behavior that causes supersonic crack speeds in mode II cracks is stiffening.}
\end{quote}

This response is accurate, concise, and clear~\cite{Buehler2008AtomisticFailure}.

The distinction between a base model and the X-LoRA model becomes even more pronounced in specialized protein analysis and design tasks. The given task is:

\begin{quote}
\small\texttt{CalculateForceHistory< M S K S N T Y R M L V L L E D D T K I N K E D E K F L K G K P G K M H E F V D E L I L P F N V D E L D E L N T W F D K F D A E I C I P Y V E I L K E S G M K > }
\end{quote}

The response from the base foundation model is:

\begin{quote}
\small\texttt{I'm sorry but the input provided is not a valid mathematical expression. It appears to be a sequence of characters that do not form a calculable function. Please provide a valid mathematical expression for me to calculate.

However, I can tell you that the sequence of characters you provided is actually a series of function and variable names that may be part of a larger program or script in a specific programming language. Without context, it's impossible to determine the intended calculation.

If you provide more information, such as the programming language and the context in which this expression appears, I may be able to provide more insight into its meaning and possible calculations. }
\end{quote}

In contrast, the X-LoRA model responds as follows, accurately completing the design task by generating the force-deformation curve linked to this protein:

\begin{quote}
\small\texttt{[0.013,0.066,0.078,0.045,0.072,0.126,0.127,0.107,0.119, 
0.074,0.106,0.098,0.103,0.128,
0.077, 0.069,0.131,0.108,0.144,0.174, 0.166,0.183,0.185,0.217,0.175,0.230,0.187,0.224, 0.267,0.301,0.315]}
\end{quote}

Figure~\ref{fig:QA_weights_sample_0} displays outcomes from this question answering task alongside the observed X-LoRA scaling weights. The two queries—one pertaining to mechanics and materials (Fig.~\ref{fig:QA_weights_sample_0}a) and the other concerning protein mechanics (Fig.~\ref{fig:QA_weights_sample_0}b)—exhibit distinct activation patterns. Notably, we observe that each case employs a complex distribution of scaling values, indicating that the X-LoRA model leverages a heterogeneous mixture of different adapters across layers. Additionally, the heatmap in Fig.~\ref{fig:QA_weights_sample_0}a reveals a wider range of activations spanning numerous experts and layers, including some regions of strong activation. Conversely, the heatmap in Fig.~\ref{fig:QA_weights_sample_0}b demonstrates a more focused activation pattern, with fewer widespread high values. We interpret this as suggesting that the gating function depicted on the right is more selective when activating experts.

In Fig.~\ref{fig:QA_weights_sample_0}, the bright line-like areas indicate paths of high activation, signifying that the scaling function selectively enhances certain experts for specific layers. The observed patterns imply that the importance is not evenly distributed across all experts and layers; instead, it tends to be quite sparse and selective. This sparsity is typical of what has been seen in other mixtures-of-experts models, where only a subset of experts is activated for each input to specialize in different aspects of the problem domain.

By comparing the activation maps across different tasks, both the sparsity and the specific patterns could be crucial for better understanding which experts the model depends on for various types of inputs or at different stages of processing within the neural network (we will later discuss how these maps evolve dynamically when solving complex tasks).

To gain deeper insight into the behavior, we now examine a broader range of questions to investigate how the X-LoRA model employs the different experts in response to diverse prompts. Our focus is on whether, and if so how, the X-LoRA model combines multiple adapters for questions that overlap across different skills or domains. Fig.~\ref{fig:QA_weights} presents results from question answering along with the observed X-LoRA scaling weights, and the corresponding queries are listed in Table~\ref{tab:table_qa}. In Table~\ref{tab:table_qa}, we provide responses from both X-LoRA and the base model, Zephyr-7B-$\beta$. The differences between the two models are notable.

To illustrate this clearly, incorrect or questionable answers are highlighted in \redhighlight{red highlight}. The comparison indicates that X-LoRA achieves substantially higher accuracy and delivers more concise responses.

As anticipated, we find complex combinations of adapters and frequent activation of multiple dominant LoRA experts. For example, Fig.~\ref{fig:QA_weights}d shows that a question about pine cone mechanics strongly activates the mechanics/materials adapter alongside the CoT/reasoning and bioinspired adapters. Although X-LoRA selects the correct letter, the answer is not fully accurate since pine cones open in dry, not wet, conditions. A subsequent question probing whether release occurs in dry or wet conditions yields the correct response.

In a protein-related question, Fig.~\ref{fig:QA_weights}f demonstrates adapter mixing in response to an inquiry about the mechanics of Bombyx mori silk, engaging a combination of bioinspired, mechanics/materials, and biology adapters. For highly specialized tasks such as protein design (Fig.~\ref{fig:QA_weights}a), we observe predominantly singular activation of one expert.

A more detailed examination of the heatmap of scalings shown in Fig.~\ref{fig:QA_weights}a reveals that, besides the clearly significant role of the protein mechanics expert, there exist several bands where particular experts receive considerable emphasis over multiple layers. This trend is especially evident for the mechanics/materials and reasoning experts. These bands are discontinuous and exhibit a fluctuating pattern, potentially indicating that these experts effectively address specific conditions or feature sets. Notably, the reasoning expert (positioned immediately to the left of the protein mechanics expert) is prominent, with several layers where it achieves the highest activation, particularly in the upper layers (approximately layers 100 to 140). This observation may suggest a crucial influence in the model’s ultimate decision-making or output, although further investigation is required to confirm this hypothesis. Additionally, the physics and biology experts show strong but more localized activations at certain layers, which could imply a specialization for processing features or representations relevant to those layers.

Moreover, Fig.~\ref{fig:QA_token_history} illustrates a token-level progression based on the aggregate scalings across all layers. This analysis presents the histories for two cases: Fig.~\ref{fig:QA_token_history}a relates to a pine cone seed release task, while Fig.~\ref{fig:QA_token_history}b pertains to a sequence of protein mechanics tasks. The protein mechanics example involves a sequence of prompt-answer interactions, beginning with two property computations followed by a generative task aimed at designing a new protein towards the end. Although expert usage remains largely consistent throughout, noticeable shifts occur in the employment of the LoRA experts during the process. In particular, the protein mechanics expert becomes highly concentrated in the final stages of generation when the new protein design is underway.

Table~\ref{tab:table_qa} presents a range of use cases spanning various domains, alongside a performance comparison with the base model alone. We elaborate on several of these cases in greater detail.

For example, we present a question that integrates aspects of protein science and biology with logical reasoning, focusing on protein expression patterns. X-LoRA provides the correct answer. The response from the base model, however, contains some errors and logical flaws when evaluating the expression of marker proteins in each cell type. The base model accurately recognizes that fibroblasts express Protein B and do not express Protein A, based on the mutual exclusivity between Protein A and Protein B. Nevertheless, it fails to explicitly mention that fibroblasts express Protein C. Since there is no rule prohibiting the co-expression of Protein C with Protein B, and fibroblasts have no restriction against Protein C, it follows logically that they also express Protein C. Regarding neurons, the base model begins correctly by stating that neurons express Protein A and not Protein B due to the mutual exclusivity rule. However, the explanation regarding the possibility of neurons expressing Protein C or both Protein C and Protein B is unclear and contains an error. Given that neurons express Protein A and not Protein B (as established), the only remaining option according to the rules is that they express Protein C. The initial analysis mistakenly suggests a scenario where neurons might express Protein B, which is incorrect. Therefore, neurons express Protein A and Protein C, not Protein B. In the case of epithelial cells, the base model accurately notes that these cells do not express Protein A but do express both Protein B and Protein C. However, the rationale provided is convoluted and includes incorrect claims about the mutual exclusivity between Protein A and Protein C, which was not indicated in the original observations. The straightforward logic is that epithelial cells do not express Protein A but must express the other two proteins, Protein B and Protein C, since the initial instructions only state that Protein A and Protein B cannot be co-expressed.

X-LoRA, on the other hand, accurately determines the marker proteins expressed by each cell type by utilizing the given observations and logical inferences. Regarding fibroblasts, X-LoRA correctly infers that they cannot express Protein A because the presence of Protein A excludes Protein B expression. Since fibroblasts do express Protein B, and given that each cell type expresses a distinct combination of three marker proteins (which was an incorrect assumption introduced, as the original statement did not specify the exact number of proteins each cell type expresses but only that the combination is unique), the conclusion that fibroblasts also express Protein C is justified based on the information available (excluding the misunderstanding about the number of proteins). For neurons, the model accurately identifies that they express Protein A and Protein C. The logic is valid, noting that Protein B cannot coexist with Protein A, and since neurons express Protein A along with another marker protein (not Protein B), they must express Protein C. Lastly, for epithelial cells, X-LoRA correctly deduces that they express Protein B and Protein C, because epithelial cells do not express Protein A and, through elimination and the fact they express two marker proteins, these proteins must be Protein B and Protein C.

Regarding the question about pine cone release under humid versus dry conditions, the base model provides a confused and incorrect response, whereas X-LoRA offers a clear and precise explanation with sound reasoning. Similar cases appear in other examples, such as the question about the unique mechanical properties of Bombyx mori silk, where the base model mistakenly claims that silk fibers conduct electricity. The base model also fails at various domain-specific tasks, such as the question about the role of hyperelasticity in maximum fracture speeds, where X-LoRA delivers a succinct and correct answer, and the base model does not respond adequately. In the query concerning the strong adhesion of gecko feet, X-LoRA rightly highlights the significance of setae and nanoscale effects, whereas the base model erroneously attributes the adhesion to special glue proteins. These and other instances detailed in the table and elsewhere in the paper strongly support the superior performance of X-LoRA compared to the base model.

Figure~\ref{fig:perf_comp_bioinspiredexam} presents a comparison of knowledge recall evaluation experiments conducted using the bio-inspired knowledge test set introduced in \cite{Luu2023BioinspiredLLM:Materials}. It is evident that X-LoRA achieves the highest performance, despite being a significantly smaller model than the BioinspiredLLM or Llama-BioLLM models (7B parameters for X-LoRA versus 13B parameters for the others). The figure also includes comparisons with other recently published large language models, such as Orca-13B and Llama-13b-chat.

Additionally, the plot compares X-LoRA to its base model, Zephyr-7B-$\beta$, demonstrating that X-LoRA delivers markedly improved performance. These results suggest that substantial performance gains can be attained with X-LoRA even when smaller base models are employed.

Figure~\ref{fig:image_synth} illustrates an application in image synthesis, where X-LoRA is used to generate a prompt intended for use in other models, such as DALL-E 3 in this case.

\begin{quote}
\small\texttt{Create a prompt that I can use for image generation that accurately describes the microstructure of a hierarchical composite.

Explicitly describe the geometry.}
\end{quote}

The comparison shown in Figure~\ref{fig:image_synth}(b)-(c) clearly reveals that the prompt generated by X-LoRA is more concise, which leads to a significantly more accurate image produced by DALL-E 3~\cite{DALLECard}. Notably, the prompt created by Zephyr-7B-$\beta$ fails to faithfully reflect the given instruction and introduces additional design elements such as fibers and nanoparticles, which were not mentioned in the original prompt that specifically targeted the microstructure of a hierarchical composite. The X-LoRA prompt is both shorter and more targeted, thereby yielding a superior result.

\subsection{Generative solutions to protein analysis}

We now turn to targeted generative and analytical tasks involving proteins. These functionalities arise in the X-LoRA model through the protein mechanics adapter and are further strengthened by integration with the model's existing knowledge in biology, bio-inspired materials, reasoning, and logical deduction. This section focuses on analyzing protein-related tasks independently, evaluating the quantitative performance of this capability.

Fig.~\ref{fig:protein_force_history_prediction} presents outcomes related to the protein task, which involves predicting force-deformation behavior from sequences. Fig.~\ref{fig:protein_force_energy_prediction} illustrates the overall accuracy in forecasting unfolding force and unfolding energy based on the sequence, comparing actual values with predicted ones (with an observed $R_2=0.85$). In general, the model demonstrates excellent performance and, as will be discussed in the following section, also performs exceptionally well in the inverse task of design and cycle consistent evaluation\cite{Lu2023GenerativePropertiesb}. Unlike previous studies that employed highly specialized GPT-like models for such tasks, our X-LoRA model integrates these specialized capabilities with a wide range of additional deep specializations achieved through various fine-tuned adaptations.  

\subsection{Protein design and analysis}

We now broaden the scope of test cases to encompass more integrative applications that engage multiple areas of capability. Specifically, the experiments described in this section demonstrate how the X-LoRA model leverages a diverse set of skills and how agentic modeling can produce more complex and comprehensive responses. We start by addressing a protein design challenge, where the model is tasked with creating a novel protein that exhibits a specified target force-deformation response. These force-deformation profiles have been examined in prior research~\cite{Ni2023ForceGen:Model} and correspond to measurements of force versus deformation when the protein is stretched at both ends, representing a classical protein unfolding experiment~\cite{Lu1998UnfoldingSimulation.}.

Fig.~\ref{fig:design_protein} illustrates the application of the generative protein task. We design a protein exhibiting a specified force-deformation response (training data and boundary conditions from the original molecular dynamics (MD) simulations of protein pulling, as detailed in~\cite{Ni2023ForceGen:Model}), and subsequently evaluate the predicted sequence to determine if it achieves the intended target properties (Fig.~\ref{fig:design_protein}a). This cycle-consistent evaluation is crucial for verifying that the inverse design capability aligns with the forward prediction.

Fig.~\ref{fig:design_protein}b presents the sequence predicted: \texttt{MSKSNTYRMLVLLEDDTKINKEDEKFLKGKPGKMHEFVDELILPFNVDELDELNTWFDKFDAEICIPYVEILKESGMK}. Additionally, Fig.~\ref{fig:design_protein}c displays a folded model of the protein generated using AlphaFold 2~\cite{Jumper2021HighlyAlphaFoldb}. To investigate the designed protein's relationship with known sequences, Fig.~\ref{fig:design_protein}d shows a Basic Local Alignment Search Tool (BLAST) Tree analysis~\cite{Altschul1990BasicTool}. Notably, the designed protein shares connections with hypothetical proteins examined in prior studies, which have not yet been identified or structurally resolved.

Leveraging agentic modeling and incorporating multiple X-LoRA models as components of a multi-agent framework, we further evaluate the predicted sequence using our X-LoRA model. This approach exploits our model’s enhanced ability not only to comprehend protein mechanics but also to integrate this knowledge with other domains. We use the sequence prediction as input and request an in-depth analysis and design recommendations for sequence variants. We initiate the task with the following prompt:

\begin{quote}
\small\texttt{I have identified this amino acid sequence: MSKSNTYRMLVLLEDDTKINKEDEKFLKGKPGKMHEFVDEL ILPFNVDELDELNTWFDKFDAEICIPYVEILKESGMK.

How can I use it?

Provide logical steps, discuss specifically the engineering of the problem.}
\end{quote}

We engage multiple agentic interactions employing adversarial strategies (see Materials and Methods for details) to encourage the model to probe more intricate and nuanced aspects of the topic. The X-LoRA model’s summarized outcome of this discussion is presented below (for the complete dialogue, refer to Supplementary Material, SI Conv Example 1):

\begin{quote}
\small\texttt{The discussion centers on the possible uses of a particular amino acid sequence in protein engineering. This sequence comprises both hydrophobic and hydrophilic residues, indicating its potential to form a stable structure and interact with other molecules. Computational approaches can be employed to predict the secondary and tertiary structures when designing a protein with this sequence; however, challenges include the complexity of protein folding and the high computational demands. Enhancements can be achieved by creating more efficient algorithms and integrating new experimental findings. It is essential to validate computational models to guarantee their accuracy and relevance for practical use. Methods like cross-validation, regularization, early stopping, ensemble techniques, and transfer learning can help avoid overfitting and underfitting. Additionally, interpreting and explaining the models can provide valuable insights for practical applications. }
\end{quote}

We also request the model to generate a summary of main points, leading to:

\begin{quote}
\small\texttt{\begin{itemize}
    \item The amino acid sequence MSKSNTYRMLVLLEDDTKINKEDEKFLKGKPGKMHEFVDELI LPFNVDELDELNTWFDKFDAEICIPYVEILKESGMK holds promise in protein engineering because of its combination of hydrophobic and hydrophilic residues. Computational techniques can predict the protein’s secondary and tertiary structures, though protein folding complexity and computational expense pose challenges.
    \item Validating computational models is vital to ensure their precision and applicability in real-world scenarios. Strategies such as cross-validation, regularization, early stopping, ensemble methods, and transfer learning help mitigate overfitting and underfitting.
    \item Understanding and explaining models also supports their effective application.
\end{itemize}
}
\end{quote}

In conclusion, the key insight is:

\begin{quote}
\small\texttt{The key insight is that the amino acid sequence MSKSNTYRMLVLLEDDTKINKEDEKFLKGKPGKMHEFVDELILPFNVDELDELNTWFDKFDAEICIPYVEILKESGMK holds promise for protein engineering because of its combination of hydrophobic and hydrophilic residues. Computational techniques can be employed to predict its secondary and tertiary structures for protein design, though challenges arise due to the complexity of protein folding and the computational resources required. It is essential to validate these computational models to ensure their reliability and practical applicability. Approaches such as cross-validation, regularization, early stopping, ensemble methods, and transfer learning are useful to mitigate overfitting and underfitting. Additionally, interpreting and explaining model outputs can provide valuable insights for real-world implementation.}
\end{quote}

We can further delve into particular sequence engineering strategies. To achieve this, we prompt the model for more targeted design recommendations aimed at enhancing the protein’s stability:

\begin{quote}
\small\texttt{How could I modify the following protein sequence to improve its stability: MSKSNTYRMLVLLEDDTKINKEDEKFLKGKPGKMHEFVDELILPFNVDELDELNTWFDKFDAEICIPYVEILKESGMK? }
\end{quote}

The summary of the dialogue is (complete conversation available in Supplementary Material, SI Conv Example 2):

\begin{quote}
\small\texttt{The discussion focuses on strategies to increase protein sequence stability by incorporating hydrophobic amino acids and disulfide bonds. The biologist explains that hydrophobic amino acids often contribute to protein stability, and the addition of disulfide bonds can further enhance this effect. The engineer requests concrete examples of hydrophobic amino acids and guidance on the number of disulfide bonds to introduce. The biologist cites instances of proteins successfully stabilized using these modifications, such as insulin, bovine pancreatic trypsin inhibitor, and green fluorescent protein. The conversation also addresses potential challenges or disadvantages of adding hydrophobic residues or disulfide bonds, including alterations in protein function, increased structural complexity, and risk of aggregation. To overcome these issues, the biologist recommends careful amino acid selection, computational analyses, and experimental validation.}
\end{quote}

The list of main points:

\begin{quote}
\small\texttt{\begin{itemize}
\item To enhance the stability of a protein sequence, one can introduce hydrophobic amino acids and disulfide bonds.

Hydrophobic residues generally contribute to greater protein stability, and disulfide bonds similarly improve structural robustness.

Examples of proteins successfully stabilized using these modifications include insulin, bovine pancreatic trypsin inhibitor, and green fluorescent protein.
\item Possible drawbacks of adding hydrophobic amino acids or disulfide bonds include alterations in protein function, increased complexity, and a higher risk of aggregation.
\item These issues can be mitigated by careful amino acid selection, computational modeling, and experimental verification.
\end{itemize}
}
\end{quote}

Ultimately, the essential conclusion is:

\begin{quote}
\small\texttt{The most critical point from the discussion is that to enhance the stability of the protein sequence MSKSNTYRMLVLLEDDTKINKEDEKFLKGKPGKMHEFVDELILPFNVDELDELNTWFDKFDAEICIPYVEILKESGMK, hydrophobic amino acids and disulfide bonds can be incorporated. This directly addresses the initial query.}
\end{quote}

From the viewpoint of protein engineering, these recommendations are sensible as they would indeed contribute to more stable designs. Concrete examples and strategies are provided that researchers can investigate further.

In a different example, three sequences are examined. First, the model is asked to assess the strength of each sequence, then to reason about the findings: identify the strongest candidate and justify why that protein is likely the most robust. Additionally, the model is requested to describe the steps for its laboratory synthesis. The three sequences (generated by manual alteration of one sequence from the test set) are:

\begin{quote}
\small\texttt{
\begin{itemize}
    \item LNTWFDKFDAEICIPYVEILKESGMK \item AITWFDKFDAEICIPYVEALKIAAMK \item AAAAAAAAKFDAAEICIIAAAAAAAAKIAAMK
\end{itemize}
}
\end{quote}

This example assesses whether the model can integrate multiple capabilities and adapt dynamically throughout the conversation. Specifically, it tests the ability to switch flexibly among different skills such as understanding protein mechanics, performing rational analysis, generating novel hypotheses, and so forth.

The conversation is outlined in Figure~\ref{fig:conversation_evolve}, which also illustrates how the scalings evolve throughout the discussion (the analysis includes both the detailed layer-wise adapter scaling and an aggregated analysis to provide an estimate of the effective utilization of the different experts). The findings clearly indicate that the protein task is initially highly relevant, progressively transitioning to a stronger emphasis on the bioinspired adapter and the biology adapter. By the conclusion of the conversation, the biology adapter becomes the most dominant. 

To evaluate the predictions and assess the quality of the responses, the three proteins involved in this task are shown in Figure~\ref{fig:protein_visuals} (proteins folded using Alpha Fold 2~\cite{Jumper2021HighlyAlphaFoldb}, via ColabFold \cite{Mirdita2022ColabFold:All}). We observe that the most stable structure in the center displays the most organized secondary structure, consistent with the idea that it produces the highest unfolding force, as forecasted by the X-LoRA model in Figure~\ref{fig:conversation_evolve}. The other two proteins are less organized; the first is mainly helical with some turns, while the last is partially unstructured and features a bend in its shape. These characteristics account for the lowest unfolding force predicted by the model.

By examining the visual depictions of the folded configurations, we verify that the structure in the center, which is the most stable, exhibits the most well-organized secondary structure. This aligns with the concept that it generates the highest unfolding force as forecasted by the model. We observe that the other two structures are less organized: the first predominantly helical with some turns, and the last one partially disordered and featuring a bend in its conformation. We suggest that these characteristics likely account for the lowest unfolding force observed, consistent with the predictions made by the X-LoRA model.

In a subsequent example, we focus on the energetic characteristics of proteins, evaluating the unfolding energy as the protein is extended from its initial folded state to a fully unfolded configuration~\cite{Ni2023ForceGen:Model}. We again analyze three sequences (the first two are the same as those previously studied, while the third is a longer construct combining the second sequence at the ends and the first sequence in the middle):

\begin{quote}
\small\texttt{\begin{itemize}
    \item LNTWFDKFDAEICIPYVEILKESGMK \item AITWFDKFDAEICIPYVEALKIAAMK \item AITWFDKFDAEICIPYVEALKIAAMKLNTWFDKFDAEICIPYVEILKESGMKAITWFDKFDAEICIPYVEALKIAAMK
\end{itemize}
}
\end{quote}

Initially, we request the model to calculate the unfolding energy for each sequence, followed by an analysis of the outcomes: based on the amino acid sequences, we ask for an explanation of why a particular protein is likely the most stable. Finally, we inquire about the potential functions that the most stable protein may perform.

Figure~\ref{fig:MJB_20} presents a transcript of a dialogue that combines different expert modules and synthesizes information from these capabilities, demonstrated here in the context of analyzing the unfolding energy of three proteins. Towards the conclusion of the discussion, the CoT/reasoning adapter becomes the dominant component as the X-LoRA model addresses a question requiring reasoning over the results to predict probable protein structural characteristics along with potential functional roles. Specifically, the mechanistic inquiry into the basis of protein stability highlights the formation of hydrophobic interactions among nonpolar amino acids and hydrogen bonds between polar amino acids. The model forecasts that several hydrophobic residues (I, F, W, Y) and hydrogen-bonding residues (D, E, K) play a role in enhancing the protein’s stability.

Additionally, the model anticipates that cysteine (C) residues in the sequence will facilitate disulfide bond formation, which further stabilizes the protein’s structure. It is noteworthy that the primary structural feature predictions, which explain the protein’s high stability, are accurately identified and can be corroborated by the structural analyses displayed in Fig.~\ref{fig:MJB_21}. This analysis confirms the existence of all these stabilizing elements. Moreover, it clearly explains why the other two proteins, which are shorter sequences with variable helical domain stability, exhibit significantly lower unfolding energies.

\subsection{Adversarial agentic modeling to connect distinct scholarly disciplines and knowledge yielding ontological knowledge graph generation}

Given that our X-LoRA model possesses extensive knowledge across multiple domains, it can be employed to explore links among diverse ideas, knowledge repositories, and fields of expertise. In one experiment, we submit two queries that are similarly structured but emphasize different aspects. Subsequently, we generate triplets that produce an ontological knowledge graph, thereby refining the responses into more organized outputs~\cite{Giesa2012CategoryDesign,Spivak2011ReoccurringAnalogies}. The resulting graph offers a cohesive understanding of the derived insights and visually represents the relationships between concepts in a comprehensible and mechanistic fashion.

We employ two adversarial agents to address the question. Agent 1 is labeled as a "critical physicist," while Agent 2 is designated as a "philosopher." Agent 1 initiates the inquiry and is responsible for probing and generating further questions in response to each reply from Agent 2. As the dialogue progresses, it delves deeper into the topic and reveals multiple facets. The question is deliberately crafted to span several domains of disciplinary knowledge to evaluate the model's capability to synthesize ideas and concepts.

The question posed is:

\begin{quote}
\small\texttt{Consider an abstract representation of a classical musical composition as a graph of interacting notes, chords, melodies, rhythms and so on.

What happens if we would apply pressure, up to a point of failure, on this abstract concept of music?

Provide logical steps, mathematical reasoning, and discuss specifically the physics of the problem.}
\end{quote}

The complete conversation is included in the Supplementary Material (see, SI Conv Example 3). For brevity, we present only a summary and the main takeaway here, which encapsulates the core aspects:

\begin{quote}
\small\texttt{The dialogue explores the effect of applying pressure to an abstract representation of a classical musical composition, treating it as a physical system comprising interacting notes, chords, melodies, and rhythms. The applied pressure can be characterized using stress, and the failure threshold is identified when the stress attains a critical value. To validate the mathematical model's predictions, a physical model of the abstract musical composition can be constructed, and experimental outcomes can be compared with the theoretical forecasts. The proposed mathematical model has limitations such as a simplified depiction of interactions, static modeling, uniform distribution assumptions, and a linear correlation between stress and pressure. To enhance the model’s precision and robustness, it should incorporate more detailed and realistic representations, dynamic factors, heterogeneous distributions, and nonlinear components.

[...]

The key insight is that applying pressure to an abstract representation of a classical musical composition can be conceptualized as altering the force field among the interacting notes, chords, melodies, and rhythms. With increasing pressure, the system’s stress rises, and beyond a critical threshold, the system will fail and collapse under its own weight. This understanding answers the original question by offering a physical rationale for how pressure influences the abstract notion of music and how it can be quantified through stress.}
\end{quote}

Although the generated text delivers valuable perspectives, constructing an ontological representation of the data can facilitate interpretable analysis~\cite{Buehler2023MechGPTModalities}. For graph generation, we merge the entirety of the two conversations and build a knowledge graph. An example of such a graph is displayed in Figure~\ref{fig:graph}. The resulting graph is partitioned into communities by applying the Girvan-Newman algorithm~\cite{Girvan2002CommunityNetworks}, yielding a total of 12 communities. The figure presents both the complete graph (Figure~\ref{fig:graph}a) and more detailed views with node labels (Figure~\ref{fig:graph}b-c).

\subsection{Creation of X-LoRA-Gemma Integrating Protein, Chemical, Bio-inspired, and Materials Mechanics Functionalities}
\label{gemma-section}

To demonstrate that our proposed method is compatible with other foundational models featuring different architectures, we trained another X-LoRA variant, this time utilizing the Gemma-7B-it model~\cite{Google/gemma-7b-itFace}. This model employs a set of four LoRA adapters, defined as follows:

\begin{enumerate}
    \item Bio-inspired materials \item Mechanics and materials science \item Protein mechanics tasks (including generative sequence-to-property and inverse design capabilities) \item Quantum mechanics-based molecular properties QM9 (incorporating generative SMILES-to-property and inverse design functions; see Table~\ref{tab:table_QM9} for definitions and a summary of all properties calculated in that dataset)~\cite{Ruddigkeit2012EnumerationGDB-17,Ramakrishnan2015ElectronicSpace}
\end{enumerate}

The X-LoRA-Gemma model is developed similarly to the original X-LoRA model but targets different capability sets, particularly adding the prediction of 12 quantum mechanical properties and the ability to design molecules satisfying these 12 quantum mechanical criteria. This new adapter is trained concurrently with the other adapters listed above, based on the same datasets as those used in the previous experiments. It is worth noting that while the majority of the paper focuses on the Zephyr-based X-LoRA model, within this section we explicitly differentiate the models by referring to the original as X-LoRA-Zephyr and the Gemma-based variant as X-LoRA-Gemma.

Figure~\ref{fig:protein_force_energy_prediction} illustrates the overall performance of X-LoRA, showing predictions of unfolding force and unfolding energy derived from the sequence ($R_2 = 0.87$), as well as the prediction accuracy for quantum mechanical molecular properties from the QM9 dataset ($R_2 = 0.96$). Both tasks exhibit excellent performance, with the protein design predictions even slightly surpassing those of the original X-LoRA model.

We proceed to showcase X-LoRA-Gemma’s abilities through a straightforward multi-agent molecular design experiment. We select a small organic molecule, \texttt{CCCC12CC1C(=O)N2} (extracted from the QM9 test set), which is a bicyclic compound containing a ketone (=O) and an amine (N) within the ring, as depicted in Figure~\ref{fig:QM9_design}(a). Our first query to the X-LoRA-Gemma model is:

\begin{quote}
\small\texttt{Here is a molecule in SMILES code: CCCC12CC1C(=O)N2

What changes would I need to make to make this molecule more useful for polymer design?}
\end{quote}

The model’s response is:

\begin{quote}
\small\texttt{To enhance the utility of this molecule for polymer design, you might consider:
\begin{enumerate}
\item Adding a reactive functional group, such as an amine or hydroxyl, to enable further modifications and crosslinking.
\item Extending the carbon backbone with more flexible chains to improve the molecule's capacity for forming polymers with diverse structures and characteristics.
\item Attaching a side chain containing a functional group like a carboxyl group to facilitate covalent bonding with other molecules and polymers.
\item Introducing a cyclic element, for example a lactam ring, to produce a polymer that is more rigid and stable.
\item Incorporating a variety of side chains to generate polymers with a broad range of properties and functionalities.
\item Varying the number of carbons in the backbone to allow for polymers with different molecular weights and traits.
\item Adding crosslinking sites to create a polymer that is more stable and less reactive.
\item Designing a repeating unit structure to synthesize a polymer with defined molecular weight and properties.
\item Combining several of these approaches to develop a versatile and functional polymer suited for multiple uses.
\end{enumerate}}
\end{quote}

Next, we concentrate on the first recommendation.

\begin{quote}
\small\texttt{How can I enhance the overall reactivity of a molecule? I intend to use this molecule for constructing a polymer material.

The molecule under consideration, represented in SMILES notation, is: CCCC12CC1C(=O)N2

I can adjust the following properties:

Dipole moment, \(\mu\) Polarizability, \(\alpha\) Highest Occupied Molecular Orbital energy, HOMO Lowest Unoccupied Molecular Orbital energy, LUMO HOMO-LUMO gap: the energy difference between HOMO and LUMO, gap Electronic spatial extent Zero-point vibrational energy Heat capacity at constant volume Internal energy at 0 K Internal energy at 298.15 K Enthalpy at 298.15 K Free energy at 298.15 K Which are the top THREE properties I should alter, and in what direction (increase or decrease) should these changes be made? }
\end{quote}

The feedback from X-LoRA-Gemma is:

\begin{quote}
\small\texttt{\begin{itemize}
    \item  Reduction in HOMO-LUMO gap \item  Increase in dipole moment \item  Enhancement of polarizability
\end{itemize}
Adjusting these characteristics can enhance the overall reactivity of the molecule, thereby improving its suitability for polymer material synthesis. }
\end{quote}

We leverage both the generative and predictive functionalities of X-LoRA-Gemma to engineer these specific changes in molecular properties.

For the optimization, we implement a two-agent system comprising two instances of X-LoRA-Gemma. One agent performs the generative task \texttt{GenerateMolecularProperties<...>} to create candidate molecules, while the other executes \texttt{CalculateMolecularProperties<...>} to evaluate the discrepancy between the target and predicted properties. This cycle continues until either the maximum iteration count is reached or the mean-squared error falls below a predefined threshold.

The outcomes of this process are illustrated in Figure~\ref{fig:QM9_design}. Panel (a) presents the design goals and their deviation from the initial molecular properties. Panel (b) displays the top 20 molecular designs produced by the algorithm, with panel (c) showing the distribution of the mean-squared errors and panel (d) plotting the error progression over successive designs. The optimal design, \texttt{CC(C)C1=CC(=O)NC=C1}, closely aligns with the target, as depicted in panel (e). Lastly, panel (f) features the optimized 3D structure of the molecule, obtained using UFF~\cite{Rappe1992UFFSimulations} for 3D conformation optimization.

To evaluate the molecule's reactivity further, we computed the Gasteiger charges~\cite{Gasteiger1980IterativeCharges}, also shown in Figure~\ref{fig:QM9_design}(f) for visualization. The molecule includes an aromatic ring with a ketone (C=O) group, which is more electrophilic than typical carbon-carbon or carbon-hydrogen bonds, rendering it prone to nucleophilic attacks. The combination of the nitrogen-containing aromatic ring and the ketone forms a pyridone ring structure. This structure exhibits greater reactivity than a standard aromatic ring because the heteroatom (nitrogen) and the ketone group provide reactive sites, such as for nucleophilic addition at the carbonyl carbon or electrophilic substitution on the ring. The nitrogen atom within the ring, as part of the heteroaromatic system, contributes to the compound's aromaticity and reactivity by modifying electron density and its distribution throughout the ring. This influences the molecule's interactions with other chemical species, making the nitrogen-containing ring a locus for potential reactions, especially electrophilic attacks facilitated by the electron-donating effect of the nitrogen atom.

The carbon atom within the ring, positioned adjacent to both oxygen and nitrogen atoms and carrying a partial positive charge (+0.25), indicates it is somewhat deficient in electrons. This occurs because oxygen, being more electronegative, draws electron density toward itself, thereby reducing the electron density on the carbon. This influence is partially offset by the nearby nitrogen, which is less electronegative than oxygen but more so than carbon. Nonetheless, the existence of the partial positive charge implies that this carbon is a probable site for nucleophilic attack, where a nucleophile (an electron-rich entity) would be attracted to the electron-deficient carbon.

The nitrogen atom, exhibiting a partial negative charge (-0.33), suggests it is relatively electron-rich compared to its normal state. This can result from its lone pair of electrons and the influence of the aromatic ring’s electron system. In aromatic heterocycles, nitrogen’s lone pair can be involved in the delocalized $\pi$-electron network, although the extent of this participation depends on the ring’s structure and the electronegativity of neighboring atoms. The partial negative charge indicates that nitrogen is more nucleophilic, meaning it has an increased tendency to donate electrons, whether through bond formation with electrophiles or protonation processes. Additional quantum mechanical studies are needed to further explore this behavior.

The carbon atom bearing the partial positive charge represents a reactive center for nucleophilic attacks. Various molecules may interact at this site via mechanisms where nucleophiles target electron-deficient carbons.

The hydrogen atom bonded to the nitrogen, having a partial positive charge of +0.17, is an excellent candidate for forming hydrogen bonds with other molecules or atoms. Hydrogen bonding is a type of dipole-dipole interaction that arises when a hydrogen atom attached to a highly electronegative atom (such as nitrogen, oxygen, or fluorine) interacts with another electronegative atom possessing a lone electron pair. This characteristic enables the molecule to act as a hydrogen bond donor, which plays a crucial role in determining the molecule’s solubility in water and other polar solvents, as well as its interactions within biological environments.

The ketone group positioned next to the nitrogen-containing aromatic ring imparts distinctive reactivity that can be harnessed in polymer synthesis~\cite{McQuarrieSimonPhysicalChemistry,CareySundbergAdvancedOrganicChemistry,Varghese2022BeyondApplications,Varghese2022BeyondApplications}. This compound is classified as a lactam, due to the cyclic amide structure formed by the ketone and nitrogen within the ring. Lactams play a crucial role in polymer chemistry, especially in the production of polyamides, a category of polymers widely utilized in various materials applications. The ketone group is capable of undergoing several chemical transformations relevant to polymer chemistry, such as nucleophilic addition, which can be exploited to elongate polymer chains or to add side chains that alter the polymer’s characteristics. The nitrogen atom in the aromatic ring can markedly influence the electronic attributes of the polymer and engage in hydrogen bonding, thereby affecting properties like melting point, solubility, and mechanical performance. N-heteroaromatic compounds are also recognized for their capacity to coordinate with metal ions, enabling potential applications in catalysis or in materials possessing electronic or photonic functionalities. Furthermore, this molecule may participate in polymerization through its amide (lactam) group; polyamides produced in this manner are valued for their durability, elasticity, and resistance to wear and chemical degradation, and are employed across diverse domains ranging from textiles (such as nylon) to automotive components~\cite{Varghese2022BeyondApplications,Varghese2022BeyondApplications}.

Alkyl substituents can enhance the hydrophobic nature of the molecule, so we expect that polymers synthesized from monomers containing these alkyl groups will tend to be more hydrophobic. This influences their solubility across various solvents as well as their interactions with water. Such properties may be beneficial for uses that demand water resistance or particular solubility traits. Additionally, the presence of alkyl groups can influence the polymer's mechanical characteristics. For example, they might serve as plasticizers within the polymer matrix, thereby improving the polymer's flexibility. This could result in materials that are more flexible or exhibit improved impact resistance, contingent on the alkyl groups' length and branching.

\subsubsection{Additional benchmarking} We present further benchmark results for the X-LoRA-Gemma model, depicted in Figure~\ref{fig:perf_MMLU}, utilizing a subset of the MMLU benchmark that concentrates on chemistry, physics, and biology (fields closely related to the domains used to train our model)~\cite{Hendrycks2020MeasuringUnderstanding}. Our findings indicate that both X-LoRA models discussed here surpass their respective base models. Specifically, the X-LoRA models outperform the base models (Zephyr model: 54.6\% for the base version compared to 56.6\% for the X-LoRA version; Gemma model: 40.6\% for the base versus 52.4\% for the X-LoRA-Gemma). The highest overall accuracy is achieved by the X-LoRA-Zephyr model, although both models demonstrate fairly similar performance.

\section{Conclusion}

Multimodal LLMs hold numerous existing and prospective applications in scientific fields, coupled with a significant demand for highly specialized expertise tailored to specific purposes. Beyond general language processing, the ability to perform computational or generative tasks, as exemplified here through protein mechanics, is crucial~\cite{Luu2023BioinspiredLLM:Materials,Buehler2023MechGPTModalities,Buehler2023MeLMProblemsb}. X-LoRA, built upon open-source language models, allows for dynamic integration of expertise and knowledge across various domains, as illustrated by multiple examples, including Figure~\ref{fig:conversation_evolve}. This study demonstrated how the X-LoRA model autonomously and adaptively evolves the application of adapters, blending certain adapters in a sophisticated manner throughout the LoRA layers. Such intricate mixing of LoRA layers is consistently observed across all examined cases, like those shown in Figures~\ref{fig:QA_weights_sample_0} and \ref{fig:QA_weights}. This phenomenon presents an intriguing challenge to gain a deeper understanding of the fundamental mechanisms from a physics viewpoint, conceptualizing a LLM as a large-scale graph-structured model exploiting complex graph topologies. Thus, the mixing of adaptation experts at deep layer scales, as executed here, modifies these mechanisms to develop effective solutions.

The introduced algorithm enhances the traditional inference approach of a single forward pass by incorporating two forward passes. This process trains the model to initially 'contemplate' the question and how it might reconfigure itself prior to generating a response. This approach embodies a basic form of 'self-awareness,' enabling the model to adjust its own architecture to optimally address a task. Consequently, despite having a relatively modest parameter count of 7 billion, the model can reason across a wide range of scientific domains (including biological materials, mathematics, physics, chemistry, logic, and mechanics), greatly boosting its potential to produce innovative solutions. Moreover, this may inform the design of alternative model architectures, such as methods that extend the reconfiguration strategy to portions or the entirety of non-adapted systems. In such cases, the scaling head could dynamically reconfigure segments of a trained model or integrate multiple models. As demonstrated with X-LoRA, these strategies can serve as powerful means to enhance or combine the capabilities of different models.

We conducted multiple performance evaluations, including a comparison with significantly larger models (Figure~\ref{fig:perf_comp_bioinspiredexam}), where X-LoRA demonstrated clear superiority. The comprehensive comparison presented in Table~\ref{tab:table_qa} indicated that X-LoRA significantly outperforms the base model across a wide range of tasks and application fields. Additional comparisons involved quantitative evaluations of numerical prediction problems, such as protein mechanics and quantum-mechanical properties of molecules, with X-LoRA achieving high accuracy levels, reflected by $R_2$ values ranging from 0.85 to 0.96. In contrast, these results surpassed other reported outcomes, for example, predicting protein unfolding force using the MeLM model~\cite{Buehler2023MeLMProblemsb}, which achieved an $R_2$ of 0.78. The ability to accurately address multiple prediction challenges, as well as inverse design and advanced reasoning tasks, was highlighted in the molecular design example (Section~\ref{gemma-section}). This case further demonstrated that X-LoRA can be effectively applied to diverse base model architectures (Gemma versus Zephyr/Mistral). Future research could extend this work by developing X-LoRA models for even larger base architectures.

A limitation of the proposed method is the need for two forward passes, which may result in increased computational overhead: one pass to compute hidden states used as input to the X-LoRA scaling head, and a second pass with $\Lambda$ applied to the adapters. However, the initial pass serves the X-LoRA scaling head and leverages the inherent capabilities of the LLM, unlike approaches that require learning the knowledge through training a larger scaling head separately. To enable efficient deployment in model-serving systems, distinct key-value caches should be maintained for both the scaling and forward passes. To enhance inference speed, we created \texttt{Mistral.rs}, a LLM serving framework that implements X-LoRA in Rust~\cite{matsakis2014rust} and incorporates various optimization techniques. Nonetheless, the key benefit of our approach is its compatibility with any model without needing to modify internal logic. We consider this a significant advantage, especially given its applicability to the extensive resources available within the Hugging Face ecosystem.

One benefit of the proposed method is that it offers an easy way to integrate with any existing LLM (for example, since the X-LoRA code has been made public, it should be straightforward to use with any model within the Hugging Face ecosystem). Another benefit is that the scaling head leverages the inherent capabilities of the LLM, learning sophisticated strategies for optimally combining the adapters to produce specific answers. During X-LoRA training, it is advantageous to include complex question-answer or conversational samples that enable the model to learn the most effective combinatorial techniques. We conducted various analyses related to protein design and mechanics. For example, Figure~\ref{fig:MJB_20} and Fig.~\ref{fig:MJB_21} showcase a stability evaluation of multiple protein sequences, incorporating comparisons, reasoning over the predicted outcomes, and a mechanistic explanation offering a fundamental chemical basis for the predicted behavior. As shown in Fig.~\ref{fig:MJB_21}, this was directly validated through structural analysis of the protein’s folded 3D conformation.

Regarding future directions, further investigation could explore a wider array of adapter experts, expanding on the initial expertise incorporated in our LoRA set. The X-LoRA model outperforms the base model alone or with only a single adapter, and it can address complex queries by utilizing the specialized capabilities contributed by the adapter set. There also appear to be interactions among the different adapters, suggested by the complex mixing patterns of scaling weights across layers. This phenomenon warrants additional study and could be compared with approaches like SLERP, which combine models using alternative methods. While we trained the X-LoRA scaling head using a subset of training data consisting of a few hundred samples from each original dataset used for adapter development, a more targeted design of a training dataset for this purpose could be implemented. For example, we observed that multiple-choice questions tend to engage reasoning-focused adapters, as they were trained on similar datasets. Extra prompting might be necessary to help the scaling head grasp a particular domain relevant to the question, potentially through specialized system messages or other techniques. Overall, creating appropriate training datasets is an area that merits deeper exploration, including methods aimed at promoting effective layer-wise scaling combinations. Our experiments tracking these scaling behaviors over extended conversations demonstrate that the model can dynamically adjust scaling strategies to optimally address specific tasks, which is a promising result.

We identified intriguing activation patterns of experts across layers, with insightful observations presented in the paper. This investigation can be extended further and has the potential to reveal significant understanding regarding the benefits of mixing components or sets of layers in models. Such exploration would be especially valuable when the X-LoRA approach is applied beyond protein mechanics and protein design, as demonstrated here. We leave this for future research.

Another promising direction involves deeper exploration of agentic modeling, as exemplified here by the adversarial interaction between two X-LoRA models. Future studies might incorporate more than two models to enrich interactions and introduce additional capabilities, thereby driving models into novel generative territories. This could also facilitate the creation of techniques that combine physics-based modeling or other validation steps, utilizing code-writing/executing agents or agents that employ function calling or alternative processing methods~\cite{Ghafarollahi2024ProtAgents:Learning}.

\section{Materials and Methods}

The X-LoRA codebase and sample applications are accessible at \url{https://github.com/EricLBuehler/xlora}. Model weights and datasets related to the X-LoRA discussed in this work, together with further examples, can be found at \url{https://huggingface.co/lamm-mit/x-lora}. The X-LoRA-Gemma model is available at \url{https://huggingface.co/lamm-mit/x-lora-gemma-7b}. The \texttt{Mistral.rs} inference engine is hosted at \url{https://github.com/EricLBuehler/mistral.rs}.

\subsection{Mathematical formulation of X-LoRA} We define $\mathbf{scalings}$ as a matrix containing coefficients at the token and layer levels for LoRA adapters. The X-LoRA $\mathbf{scaling}$ $\mathbf{head}$ receives hidden states from the base model and applies fully connected linear layers to predict scalings, followed by a softmax operation over the final dimension to guarantee that scalings sum to one across all experts. Between the linear layers in the X-LoRA scaling head, ReLU activation functions and dropout layers are interspersed. The size of hidden layers can be adjusted to either successfully reduce from the hidden dimension to the scaling dimension or to implement a widening-then-narrowing design analogous to the feed-forward layer in transformer blocks.

The scalings form a matrix $\Lambda \in \mathbb{R}^{s \times l \times n}$, where $s$, $l$, and $n$ stand for sequence length, number of layers, and number of adapters, respectively. For each layer, the corresponding scalings are extracted as a matrix $\lambda \in \mathbb{R}^{s \times n}$. These scalings $\lambda$ are then applied to each LoRA adapter by broadcast multiplying the respective X-LoRA scaling $\lambda_i \in \mathbb{R}^{1 \times 1 \times s}$ to the inputs of the decomposition matrices $A$ and $B$.

With $W_0 \in \mathbb{R}^{d \times k}$, $B_i \in \mathbb{R}^{d \times r_i}$, and $A_i \in \mathbb{R}^{r_i \times d}$, where the rank $r_i \ll \min(d,k)$ is specific to each adapter $i$, the forward computation of an X-LoRA layer can be represented by the following equation:

\begin{equation}
    h = W_0x + \sum_{i=1}^{n} B_iA_i(x \cdot  \lambda_i) \alpha_i
\end{equation}

Representations of the scalings are shown in the main text, visualized either through bar plots or via the \texttt{contourf} function in Matplotlib~\cite{MatplotlibDocumentation}.

\subsection{X-LoRA implementation algorithm and code} The implementation selected here prioritizes user-friendliness and the flexibility to easily adapt the method to different models, particularly those within the Hugging Face ecosystem. The X-LoRA algorithm functions by:

\begin{enumerate}
    \item Keeping the base model and adapters fixed, followed by fine-tuning for performance.
    \item Combining LoRA adapters based on the predictions made by the X-LoRA scaling head.
\end{enumerate}

All the code is developed using Python and PyTorch~\cite{Paszke2019PyTorchLibrary}. 

\subsubsection{Dual forward pass strategy for self-aware inference} Since X-LoRA needs to compute hidden states to estimate the scalings, we adopt a dual forward pass method. This approach employs an initial pass where the model evaluates the task to determine how to reconfigure itself, followed by the actual inference pass. This design introduces a form of `self-awareness' that enables X-LoRA to allocate extra computation time for reflection rather than immediate reaction.

We refer to the first pass as the scaling pass and the second as the forward pass. During the scaling pass, we set the scalings $\Lambda$ to zero. We have also tested setting the scalings $\Lambda$ to ${n}^{-1}$ and observed comparable outcomes, suggesting stable behavior (this parameter can be adjusted in the X-LoRA package configuration). To facilitate easy application of X-LoRA across various models, we implement a forward pass hook. This mechanism is crucial for coordinating the forward and scaling passes and is the principal feature enabling compatibility with all models.

\begin{enumerate}
    \item $\mathbf{Scaling}$ $\mathbf{pass}$: Performed by running the base model with adapters fixed at constant values. The hidden states produced here are used by the X-LoRA scaling head to compute $\Lambda$.
    \item $\mathbf{Forward}$ $\mathbf{pass}$: Utilizes $\Lambda$ to blend the adapters during the model’s forward pass. The resulting logits constitute the final output of the X-LoRA model.
\end{enumerate}

As a cautionary note, sharing the same key-value cache across both the scaling pass and forward pass can disrupt the correct input to the scaling head. This occurs because the key-value cache retains information and accumulates across the two passes. Therefore, we disable the key-value cache during both training and inference, and manage this appropriately in \texttt{Mistral.rs}.

\subsubsection{Freezing weights and enabling segments of the model for training} The base model and adapters remain frozen ({\it i.e.}, their weights are not updated during training), making the X-LoRA scaling head the sole trainable component within an X-LoRA model. The provided code allows for the option to unfreeze all adapters alongside the X-LoRA scaling head. Additionally, it is possible to train the entire model jointly—comprising the X-LoRA scaling head, adapters, and the base foundation model—potentially applying different learning rates. Moreover, future developments based on these concepts might enable creating scaling functions between two or more foundation models, facilitating a traceable realization of SLERP-style modeling, which represents a direction for future research.

\subsubsection{X-LoRA layers} An X-LoRA layer functions by scaling and subsequently combining the outputs from multiple adapters. Each X-LoRA layer encapsulates the original LoRA layer, thereby retaining access to the respective adapters associated with that layer.

To seamlessly incorporate X-LoRA into any model, the algorithm replaces the original LoRA layers with newly created X-LoRA layers and integrates their forward pass operations into the LoRA layer. This approach exploits features of Python's object system to enable layer replacement.

In summary, during the forward pass, each X-LoRA layer performs the following operations (refer to Figure \ref{fig:general_arch}):

\begin{enumerate}
    \item $\mathbf{Scaling}$: Scale the input to each LoRA adapter by the corresponding scaling factor $\lambda_i$.
    \item $\mathbf{Forward}$: Execute the forward computation of the encapsulated LoRA layer.
\end{enumerate}

\subsection{Efficient Inference with \texttt{Mistral.rs} Implemented in Rust} To enhance inference efficiency, we created \texttt{Mistral.rs}, a large language model serving framework written in Rust~\cite{matsakis2014rust}, which supports X-LoRA and incorporates various optimizations aimed at boosting performance. The implemented techniques include, but are not limited to, the following:

\begin{itemize}
    \item LoRA adapter weight stacking: When all $A_i$ and $B_i$ matrices across the LoRA adapters share the same dimensions, we stack these weight matrices $A_i$ and $B_i$ along the first axis.

Mathematically, given $W_0 \in \mathbb{R}^{d \times k}$, for adapter matrices $A_i \in \mathbb{R}^{r \times d}$ and $B_i \in \mathbb{R}^{d \times r}$ with $r_i \ll {\rm min} (d,k)$, we concatenate the matrices so that $A\in \mathbb{R}^{i \times r \times d}$ and $B\in \mathbb{R}^{i \times d \times r}$. The $\alpha_i$ parameters are also applied during initialization to the stacked matrices.

To implement scalings, we rearrange the dimensions so that $\lambda \in \mathbb{R} ^ {i \times 1 \times 1}$. These are then applied via broadcast multiplication.

\item Non-granular scalings To minimize how frequently the scaling operation is executed, we provide an option for non-granular scalings. This approach caches the scalings after $n$ completed runs, leveraging the KV cache to keep the sequence length fixed at one for the remainder of the request. This method significantly lowers latency by decreasing the number of calls made to the base model.

\item Fused Compute Unified Device Architecture (CUDA) kernels To enhance the forward pass efficiency, we employ fused CUDA kernels. These fused kernels allow complex operations—such as Rotary Embedding (\cite{su2021roformer}), which would normally run sequentially—to be combined into a single kernel executable in parallel on a GPU. Our implementation is based on the vLLM approach (\cite{kwon2023efficient}).

\item Quantization The inference engine supports quantization, enabling the use of a quantized base model. Preliminary experiments indicate that quantization down to 8 bits performs well even for models originally trained using \texttt{bfloat16}, for example.

\end{itemize}

\subsection{Datasets}

Each adapter is trained on task-specific datasets, summarized in Table~\ref{tab:table_loradata} along with references to the sources and original publications.

\subsection{Training strategy} Training proceeds in multiple stages. We start with the Zephyr-7B-$\beta$ model~\cite{HuggingFaceH4/zephyr-7b-betaFace}, which is based on the Mistral-7B model\cite{Jiang2023Mistral7B}, and train a series of adapters using the datasets described above. The Zephyr-7B-$\beta$ tokenizer is employed.

The first eight adapters are trained on question-answer pairs as provided in the datasets (refer to Table~\ref{tab:table_loradata}).

The protein mechanics tasks are trained similarly but utilize specific forward and inverse instruction sets. The bidirectional instruction tasks include:

\begin{quote}
\tiny {\texttt{CalculateForce< G E E C D C G S P ....> [0.310]\\
CalculateEnergy< G E E C D C G S P ....> [0.220]\\ 
CalculateForceHistory< G E E C D C G S P ....> [0.004,0.034,0.125,0.142,0.159,0.102,0.079,0.073,0.131, ...]\\
GenerateForce<0.310>\,[ G E E C D C G S P ....]\\
GenerateEnergy<0.220>\,[ G E E C D C G S P ....]\\
GenerateForceHistory<0.004,0.034,0.125,0.142,0.159,0.102,0.079,0.073,0.131, ...>\,[ G E E C D C G S P ....] }}
\end{quote}

These correspond to three distinct tasks: determining the maximum unfolding force of a protein, calculating the unfolding energy, and obtaining the force-deformation history. Each of these forward tasks is paired with a corresponding inverse task of the same type, but producing a protein sequence as output, resulting in a total of six tasks.

The rank for each LoRA adapter is set to $r=16$ with a scaling factor of $\alpha=16$. The adapters are trained on five specific target modules: \texttt{v\_proj}, \texttt{k\_proj}, \texttt{q\_proj}, \texttt{gate\_proj}, and \texttt{o\_proj}, using a dropout rate of 0.05. Typically, LoRA adapters are trained with low ranks, ranging from 4 to 32, because it has been demonstrated that increasing the rank does not usually lead to performance gains but does increase computational cost. Lower ranks yield more compact models with fewer parameters~\cite{hu2021lora}. The selection of $r$, $\alpha$, and other hyperparameters follows common practices in the literature known to perform well. Additionally, previous research by the author has explored various configurations and concluded that the values chosen here represent a reasonable compromise~\cite{Luu2023BioinspiredLLM:Materials}.

The X-LoRA scaling head is trained using approximately 2,000 samples, which are drawn as a subset from the original dataset employed for training the adapters (we also supplemented the training data with additional GPT-4-generated multiple-choice question samples). We utilize a \texttt{paged\_adamw\_8bit} optimizer with a gradient clipping threshold of 0.3, a learning rate of $2\times 10^{-4}$ with a warmup phase, and four gradient accumulation steps, using a batch size of one. The X-LoRA model is trained for about 10,000 steps, achieving a loss of approximately 0.28 by the end of training.

Given that the base Mistral architecture used here as the foundation model consists of 32 layers, and we apply adapters to 5 of each transformer layer, there are a total of 160 LoRA layers, each replaced by an X-LoRA layer. The X-LoRA scaling head predicts scaling values for each of these 160 LoRA layers (with nine LoRA experts, this leads to a total of $160\times 9=1440$ values of $\lambda_i$ calculated per token). We employ a single linear layer with ReLU activation functions and a dropout rate of $p=0.2$, comprising around 5 million parameters. The implementation supports flexible architectures, allowing the construction and exploration of more complex deep classifier heads for other use cases.

Both LoRA adapters and the X-LoRA scaling head are trained via next-token prediction using cross-entropy loss (we compute the probability assigned by the model to the correct next token in a sequence given the preceding tokens; the training objective is to maximize the likelihood of the correct next token as indicated in the training data). We employ token shifting, where the sequence is shifted by one token to generate the target sequence that the model aims to predict. For instance, if the input sequence is ``Proteins are fundamental to", the corresponding target sequence for training becomes ``are fundamental to biology". Thus, the model endeavors to forecast the subsequent token in the sequence.

As illustrated in Figure~\ref{fig:hierarchical_concept}, during the training of the X-LoRA scaling head, only the scaling head parameters are updated while all other parameters remain fixed. The standard Zephyr tokenizer and prompt template are utilized.

\subsection{X-LoRA-Gemma model} The X-LoRA-Gemma model is constructed similarly to the X-LoRA model described above but is based on the Gemma-7B-it model~\cite{Google/gemma-7b-itFace}, incorporating four adapters (refer to the main text for further information). These four adapters are trained on datasets covering mechanics and materials, protein mechanics, bioinspired materials, along with an additional dataset of quantum-mechanics based molecular properties QM9 (see Table~\ref{tab:table_QM9})~\cite{Ruddigkeit2012EnumerationGDB-17,Ramakrishnan2015ElectronicSpace}.

Each LoRA adapter has rank $r=16$ and scaling factor $\alpha=16$. Training involves the following seven target modules: \texttt{q\_proj}, \texttt{o\_proj}, \texttt{k\_proj}, \texttt{v\_proj}, \texttt{gate\_proj}, \texttt{up\_proj}, and \texttt{down\_proj}, with a dropout rate of 0.05. The chosen rank for each LoRA adapter is $r=16$ with $\alpha=16$. Given four adapters targeting seven modules each, and considering the Gemma-7B base model comprises 28 layers, this results in a total of $28\times 28=784$ $\lambda_i$ values computed per token.

The first two adapters (bioinspired and mechanics of materials) integrated into this model are trained using question-answer pairs sourced directly from their respective datasets (refer to Table~\ref{tab:table_loradata}).

The protein mechanics adapters are trained similarly to the previously described approach, utilizing bidirectional forward and inverse instruction sets. Additionally, an adapter is trained using the QM9 dataset, focusing on the following two tasks:

\begin{quote}
\tiny {\texttt{CalculateMolecularProperties<CC(C)N1CC1C=O> [0.098,0.358,0.581,0.395,0.309,0.330,0.570,0.649,0.519,0.519,0.519,0.518]\\
GenerateMolecularProperties<0.098,0.358,0.581,0.395,0.309,0.330,0.570,0.649,0.519,0.519,0.519,0.518> [CC(C)N1CC1C=O]\\ 
...
}}
\end{quote}

These encompass a total of two tasks: calculating molecular properties (the forward task) and the inverse task of generating a molecule with specified target properties. All 12 properties included in the QM9 dataset are either predicted or designed for, respectively (refer to Table~\ref{tab:table_QM9}).  
The model incorporates a total of four adapters.  

The X-LoRA scaling head in this model consists of three layers, with an intermediate layer size of $3072\times {FF}_{\rm fac}=12,288$, where the feed-forward multiplier ${FF}_{\rm fac}=4$. Linear layers with ReLU activation functions and a dropout rate of $p=0.2$ are used, amounting to approximately 47M parameters.  

The X-LoRA-Gemma scaling head is trained using about 17,000 samples. These samples were selected as a subset from the training dataset employed for adapter training. Notably, although the X-LoRA-Gemma model contained only four adapters, samples were drawn from the entire training set of the model mentioned above, supplemented by samples from the QM9 dataset. This was done to explore the possibility that the X-LoRA model might acquire additional capabilities during the scaling head training phase. Given that the loss converges well to 0.58 by the end of training, along with outstanding performance even on tasks for which the model was not explicitly trained (such as chemistry, exemplified by the molecule design case discussed in the paper), we believe this approach generally yields strong results.

We utilized a \texttt{paged\_adamw\_8bit} optimizer featuring gradient clipping set at 0.3, a learning rate of $1\times 10^{-4}$ with a warmup phase, and four gradient accumulation steps, operating with a batch size of one. The X-LoRA-Gemma model was trained for approximately 62,000 steps.

The training employed the original Gemma tokenizer and prompt template.

\subsection{Adversarial agentic modeling} We adopt an adversarial agentic approach by deploying two X-LoRA agents. One agent specializes in posing questions: it initiates the first question and continuously formulates further questions to uncover challenges and flaws in the answers. The second agent is responsible for responding to these inquiries, consistently providing answers. This methodology is implemented using \{Guidance\}~\cite{Guidance-ai/guidance:Models.}.

Regarding the example connecting music and mechanics discussed in the text, the roles of the two agents are defined as follows. First, the physicist and questioner:

\begin{quote}
\small\texttt{You are a critical physicist engaged in a discussion from a physics viewpoint. Keep your responses concise and always challenge statements in a provocative manner.

As a creative thinker, you incorporate ideas from other disciplines and push boundaries.}
\end{quote}

Second, a philosopher (in some instances, this role is adapted to represent a biology expert):

\begin{quote}
\small\texttt{You are a philosopher with expertise in biology, chemistry, and mathematics. You participate in the discussion.

Keep your replies brief, accurate, and creative. You generate excellent ideas and novel lines of thought, always maintaining logical consistency.}
\end{quote}

The question asker is explicitly directed to consider the question and preceding answers to formulate new insightful questions. This is implemented via a prompt structured as follows:

\begin{quote}
\small\texttt{Examine the following question and answer.

\#\#\# Question: <question>

\#\#\# Response: <response>

\#\#\# Instruction: Provide a SINGLE follow-up question that critically examines the response.

Do NOT provide an answer or commentary on the question at this point.

The single question is: }
\end{quote}

After the dialogue reaches $N$ turns, the model’s task is to 1) summarize the primary insights covered, 2) enumerate the significant insights as bullet points, and 3) pinpoint the most crucial conclusion.

The unprocessed conversation text serves as input for further analysis through graph generation.

\subsection{Knowledge graph generation} We employ Zephyr-7B-$\beta$ to extract triplets from the text, following the methodology described in~\cite{Rahulnyk/knowledge_graph:QnA} with added features inspired by the Llama Index graph generation technique~\cite{Liu_LlamaIndex_2022}. Graphs are rendered using NetworX\cite{Networkx/networkx:Python} and Pyvis~\cite{WestHealth/pyvis:Graphs.}. One essential directive, as recommended in~\cite{Liu_LlamaIndex_2022}, is:

\begin{quote}
\small\texttt{Your task is to extract the ontology of terms mentioned in the given context. These terms should represent the key concepts as per the context.

Present your output as a JSON list. Each list item contains a pair of terms and the relationship between them, for example: [...] }
\end{quote}

We direct the construction of triplets of nodes and edge features as follows, consistent with the method in~\cite{Liu_LlamaIndex_2022}:

\begin{quote}
\small\texttt{\begin{itemize}
            \item "node\_1": "A concept extracted from the ontology"
            \item "node\_2": "A related concept extracted from the ontology"
            \item "edge": "A concise description of the relationship between node\_1 and node\_2"
        \end{itemize}
        }
\end{quote}

We also supply the model with an example of ontological analysis, such as:

\begin{quote}
\small\texttt{
       Context: ```Spider silk is a strong natural fiber used to catch prey in a web.``` }
\end{quote}

Then, example triplets are provided, for instance:

\begin{quote}
\small\texttt{\begin{itemize}
            \item "node\_1": "spider silk"
            \item "node\_2": "fiber"
            \item "edge": "is"
        \end{itemize}
        }
\end{quote}

The generated graph data is compiled and subsequently partitioned into communities using the Girvan-Newman algorithm~\cite{Girvan2002CommunityNetworks}.

\subsection{Visualization of molecular structures}

PyMOL~\cite{Schrodinger2015The1.8} is utilized to visualize and analyze predicted protein structures. Various scripts and functionalities within the software are employed to detect specific protein features, such as secondary structures, hydrophobic/hydrophilic regions, disulfide bonds, and hydrogen bonds.

\section*{Data Availability Statement} The code and data supporting the conclusions of this work are publicly accessible at \url{https://github.com/EricLBuehler/xlora} and \url{https://huggingface.co/lamm-mit/x-lora}. Additional datasets used in this study are cited (refer to Table~\ref{tab:table_loradata}). 

\end{document}